\section{Existence: the canonical construction}

Using the reformulation established in the previous section, we present a simplicial cup-\(i\) construction satisfying all the axioms.
It plays a central role in the proof of our main theorem in \cref{s:proof}, and in \cref{ss:original} we prove that it coincides with Steenrod's original construction.
This contruction appeared first in \cite{medina2023fast_sq} and was used in \cite{medina2022per_st} to compute persistent Steenrod squares of real-world data.

\subsection{Canonical construction}

For any $U \in \rP^n_{n-i}$ the \textbf{index function} of $U$ is given by
\[
\begin{split}
	\ind_U \colon U &\to \F \\
	u_j &\mapsto u_j + j \mod 2.
\end{split}
\]
For any $U \in \rP^n_{n-i}$ and $\varepsilon \in \Ftwo \cong \{0,1\}$ we write $U^\varepsilon$ instead of $\ind_U^{-1}(\varepsilon)$.
Inspecting \cref{ss:canonical} we see that the canonical \mbox{cup-$i$} construction is defined by the elements
\[
\canonical_i[n] =
\sum_{\mathclap{\quad U \in \rP^n_{n-i}}} \ U^0 \ot U^1.
\]

\begin{theorem}\label{t:existence}
	The canonical \mbox{cup-$i$} construction is simplicial, non-zero, irreducible, and free.
\end{theorem}

\begin{proof}
	It is clearly non-zero and a straightforward application of \cref{l:dual_axioms} establishes its irreducibility and freeness.
	To see it is simplicial we consider \cref{eq:codegeneracies} and an arbitrary basis element $U^0 \ot U^1$ appearing as a summand in $\canonical_i[n]$.
	Let $U = U^0 \union U^1$.
	Based on \cref{l:simplicial_from_semisimplicial}, we need to check that for any $j \in \set{0,\dots,n-1}$ we have
	\begin{equation}\label{eq:checking_simplicial}
		\cP(\sigma_j)(U^0) \ot \cP(\sigma_j)(U^1) = 0.
	\end{equation}
	If $j \notin U$ or $j+1 \notin U$ either $\cP(\sigma_j)(U^0) = 0$ or $\cP(\sigma_j)(U^1) = 0$, so Identity \eqref{eq:checking_simplicial} holds.
	Let us assume $j,j+1 \in U$.
	Since $\ind_U(j) = \ind_U(j+1)$ then either $j,j+1 \in U^0$ or $j,j+1 \in U^1$.
	In both cases we also have \eqref{eq:checking_simplicial}.
\end{proof}

\subsection{Special cases}\label{ss:canonical_special_cases}

We close by recording two important instances of the canonical construction.
For any $n \in \N$,
\[
\canonical_n[n] = \emptyset \ot \emptyset,
\]
and, for $n > 0$,
\[
\canonical_{n-1} [n] \ = \
\sum_{\mathclap{\substack{\quad u \in \{0,\dots,n\} \\ u \ \mathrm{odd}}}} \,\{u\} \ot \emptyset \ + \
\sum_{\mathclap{\substack{\quad u \in \{0,\dots,n\} \\ u \ \mathrm{even}}}} \,\emptyset \ot \{u\}.
\]

%\subsection{Canonical cup-\textit{i} construction}\label{ss:canonical}
%
%To prove our main result we will use the \mbox{cup-$i$} construction introduced in \cite{medina2023fast_sq}.
%Before we recall its definition, let us introduce the following notation:
%For any $n$-simplex $x$ and set $U = \{u_1 < \dots < u_r\} \subseteq \set[big]{0,\dots,n}$ we write $d_U(x)$ for the composition of face maps $d_{u_1}\!\! \dotsm d_{u_r}(x)$ with $d_{\emptyset}(x) = x$.
%
%Let $X$ be a semi-simplicial set, $x$ an $n$-simplex of $X$, and $\alpha, \beta$ two cochains in $\cochains(X)$, then, by definition,
%\[
%(\alpha \cup_i \beta)(x) =
%\begin{cases}
%	(\alpha \ot \beta) \sum d_{U^0}(x) \ot d_{U^1}(x) &
%	\text{if } i \in \{0, \dots, n\}, \\
%	\hfil 0 &
%	\text{otherwise},
%\end{cases}
%\]
%where the sum is taken over all $U = \{u_1 < \cdots < u_{n-i}\} \subseteq \{0, \dots, n\}$ and
%\[
%U^0 = \{u_j \in U \mid u_j + j \equiv 0 \text{ mod } 2\}, \qquad
%U^1 = \{u_j \in U \mid u_j + j \equiv 1 \text{ mod } 2\}.
%\]
%As an example we mention that, if $i = 0$, the above formula gives
%\[
%(\alpha \cup_0 \beta)(x) =
%\sum_{j=0}^n \alpha \big(d_{j+1} \cdots d_{n}(x)\big) \cdot \beta \big(d_{0} \cdots d_{j-1}(x)\big),
%\]
%the well-known Alexander--Whitney product.
%
%The fact that these formulas define a semi-simplicial \mbox{cup-$i$} construction is proved in \cite{medina2023fast_sq}.
%They also underlie the algorithm for persistent Steenrod invariants introduced in \cite{medina2022per_st} and implemented in \texttt{steenroder}.
%We call this the \textbf{canonical \mbox{cup-$i$} construction}.